\section{Introduction}
A finite-state realization of a numerical experiment must preserve its
response to every continuation of the external command word. A low-rank
factorization at each cut is necessary, but need not admit transition
rows between the cuts. This distinction persists even when the target
is orthogonal, the state dimension is fixed, and every one-direction
word-distortion profile vanishes. We study the resulting compatibility
problem, with the number of available persistent labels as resource.

Fix orthogonal matrices $U_a$ on $\R^D$, a finite seed set $X$, a signal
$0<\rho<1$, and a terminal coordinate query $j$. A word
$w=a_1\cdots a_N$ prescribes the binary law
\begin{equation}\label{eq:target}
 \Pp(B=b\mid x,w,j)=\frac{1+b\rho\langle e_j,U_wx\rangle}{2},
 \qquad U_w=U_{a_N}\cdots U_{a_1},\quad b\in\{-1,1\}.
\end{equation}
The command word is supplied externally. There is no distinguished
probability measure on words. A realization may depend on the horizon
and epoch, but all persistent information, including any random selector,
must fit in its current label. Its rows may be real and are sampled
atomically. This is a nonuniform positive-realization invariant, not
ordinary workspace or program length. Definition~\ref{def:machine}
specifies it without an implicit uncharged register.

The principal result identifies a regime in which the observed response
space itself is a complete state geometry. Take $X=\{\pm e_i\}_{i=1}^D$
and allow all coordinate queries. If the identity is a command, every
normalized cut Hankel matrix has rank $D+1$. At a cut with exactly that
many labels, equality of ranks forces every label's entire future row
into the target affine response space. The row normalization then turns
a linear-space equality into an affine one. Thus no hidden future
coordinate remains, even in a machine with wider cuts before and after
it. Exact realization at rank-tight cuts is equivalent to the transport
of enclosing simplices by all intervening words
(Theorem~\ref{thm:simplex50}). This conclusion is not asserted for a
larger factorization, where hidden future directions can persist.

The simplex description has a quantitative consequence not available
from separate cut ranks. Put $c_D=2^{-D}\binom{2D}{D}$. When $I$ and
$-I$ are commands, any two consecutive rank-tight cuts force the later
simplex to contain both signs of the earlier one. The elementary
simplex difference-body identity gives a volume factor $c_D>1$ for
$D\ge2$. Therefore every exact machine has at most
\begin{equation}\label{eq:intro-budget50}
 1+\frac{\log(D!\rho^{-D})}{\log c_D}
\end{equation}
rank-tight cuts, irrespective of the widths between them. This is an
obstruction to an entire realization class, rather than a verification
of a proposed set of rows. In dimension two it gives an exact frontier:
for any signed-permutation alphabet containing $I,-I$,
\[
 (3/2)^N>2\rho^{-2}\quad\Longrightarrow\quad W_{N,0}=4.
\]
The upper bound is a common four-label permutation realization. Such
alphabets may be noncommutative. At each fixed qualifying horizon the
frontier also persists at sufficiently small positive error. We give
no numerical lower bound for that last error interval.

An all-horizon statement strengthens this finite obstruction. Exact
rank-tight realization at every horizon, even with complete nonuniform
redesign, is equivalent to one invariant full-dimensional simplex
containing $\rho\conv X$ and lying in the coordinate cube. The same
machine then works at every horizon, its command rows are permutations,
and the generated orthogonal group embeds in $\mathfrak S_{D+1}$
(Theorem~\ref{thm:allhorizon50}). In the plane and at signal at most
$1/2$, this characterizes all-horizon three-state feasibility by
containment in the symmetry group of an equilateral triangle.

These results also connect the finite theory to the word profiles
used in the arithmetic part of the article. For the alphabet
$\{I,-I\}$, every one-direction orbit has only two points. Hence
$\Gamma_{\calA}(k,B)=0$ and $\mathfrak r_{\calA}(k)=0$ for every
$k\ge2$. Nevertheless the common enclosure profile satisfies
\[
 \lambda_{\calA}(D+1,a)=D\qquad(0<a\le1/D).
\]
Theorem~\ref{thm:separation50} proves this equality by barycentric
coordinates. Thus the distortion--dilation inequality is genuinely
strict, while compatibility has a positive occupation charge. The
arithmetic results concern regimes where distortion and enclosure
scales can match; they are not consequences of a putative equality
of these two invariants in general. This distinction supplies the
conceptual link between the finite and asymptotic parts of the paper.

At two states a different geometry is available. For paired seeds and
binary responses, the odd part of a legal machine is a bounded
rank-one tensor. The complete sign--magnitude alternative retained in
Theorems~\ref{thm:duality49} and~\ref{thm:algebraic49} treats explicitly
listed tables, including zeros, by support reduction and one logarithmic
linear system per sign chamber. Integer multiplicative alternatives
exclude all machines in that class. Proposition~\ref{prop:complexity50}
prices the chamber count, sparse rays, boundary polynomials and
algebraic minimizing rows; finiteness is not confused with polynomial
complexity in a succinct horizon.

We give that tensor theory a genuinely multi-epoch orthogonal
application. Two rational noncommuting permutation matrices in dimension
nine, rational unit seeds, three command epochs and one selected
coordinate have exact two-state error
\[
 \frac{100\rho}{357}t,\qquad
 500t^3-375t^2+540t-124=0,\qquad 0.26<t<0.261.
\]
A five-entry Segre circuit proves the lower bound, and legal stochastic
rows attain it; these rows are algebraic when $\rho$ is algebraic. Every separate flattening of the positive
three-mode response admits strictly smaller rank-one error. The
obstruction therefore requires compatibility across modes, not a
single matrix relaxation. This experiment has only its designated
query; it is not an assertion about all nine coordinate queries.
The earlier arbitrary-table cubic and commuting orthogonal examples
are preserved with their distinct scopes.

Invariant positive cones, tensor exponent relations, monomial
linearization, Farkas' lemma and simplex difference-body identities are
classical tools. We distinguish their use from the realization results
proved here. The comparison in Section~\ref{sec:complexity50} includes
positive realization \cite{BF04}, matrix infinity-norm approximation
\cite{GS19}, rank-one tensor completion \cite{KKKR17} and geometric
programming \cite{BKVH07}. In particular, the cubic circuit itself is
not claimed as new. The new conclusions concern rank-tight causal
geometry, its occupation and all-horizon consequences, the exact
profile separation, and the orthogonal realization of a sharp
multi-epoch interval obstruction.

The subsequent sections retain the full two-state, positive-Hankel,
Gram, moment, arithmetic and finite-bit theorems and their proofs.
The finite-bit implementation has its own resource accounting; it
does not make exact algebraic sampling free in a uniform model.
The source-bound checks certify specified finite identities and
build reproducibility, not the universal proofs or independent
priority. The separate analytic dependencies of the wider Theta
program are neither assumed nor declared closed by the present
finite-realization results.
